\chapter{总结与展望}

\section{工作总结}

本项目围绕"大数据分析与可视化"这一主题，设计并实现了一个基于 Spark 与 Flask 的大数据分析与可视化系统。通过本项目的实践，主要完成了以下工作：

\begin{enumerate}
  \item 成功搭建了基于 Docker 的 Hadoop + Spark 大数据开发环境；
  \item 实现了基于 Spark SQL 的数据 ETL 处理和多维度分析流水线；
  \item 构建了基于 Flask 的 RESTful API 服务，提供标准化的数据访问接口；
  \item 开发了基于 React + ECharts 的可视化仪表盘，实现了分析结果的直观展示；
  \item 形成了完整的课程大作业技术文档，为后续学习和工作积累了经验。
\end{enumerate}

\section{不足与展望}

尽管本系统实现了基本的大数据分析与可视化功能，但仍存在以下不足：

\begin{itemize}
  \item \textbf{实时性不足}：当前系统主要面向离线批处理场景，未来可引入 Spark Streaming 实现实时数据处理；
  \item \textbf{分析维度有限}：当前仅实现了基础的数据分析功能，未来可引入更多高级分析算法，如聚类分析、关联规则挖掘等；
  \item \textbf{交互性有待提升}：前端可视化交互功能相对简单，未来可增加数据下钻、联动筛选等交互功能；
  \item \textbf{自动化程度不足}：数据采集、处理任务调度等流程尚未完全自动化，未来可引入 Airflow 等调度框架实现工作流自动化。
\end{itemize}

\section{心得体会}

通过本次课程大作业的实践，加深了对大数据技术栈的理解，特别是对 Hadoop 分布式存储、Spark 分布式计算、Flask Web 服务、前端可视化等技术的综合运用有了更为深入的认识。同时，在项目开发过程中也锻炼了问题分析和解决能力，为今后从事大数据相关领域的工作打下了基础。
